% Cote 119: offprints, reports, courses, CV and letters (typed and printed part): archive cote 119, site page 165 (PDF page 166).
% typescript. Draft compared: none.
% AI-transcribed by Claude Opus 5.5 (claude-opus-5-5, adaptive thinking, effort high), Message Batches API, batch r119_full (msgbatch_01DwZNm5Y7YsKBWyAN4sFs8w), 2026-11-13; not yet checked by eye.
\scanpage{119/165}
\mspage{- 4 -}

\uline{dérivée $t'(u)$ est partout $> 0$} et telle que $t(a_1) = a$, $t(b_1) = b$
$u \to \gamma\,(t(u))$ définit un chemin différentiable $\gamma_1$. La forme différentielle
$f_1(u).du$, transformée de $\omega$ par $u \to \gamma(t(u))$, est égale à $f(t(u))\ t'(u)$ du
On a donc $\int_\gamma \omega = \int_{\gamma_1} \omega$ . D'où :

\uline{Proposition 2.3}

L'intégrale curviligne $\int_\gamma \omega$ ne change pas de valeur si on remplace
le chemin différentiable $\gamma$ par un autre déduit de $\gamma$ par changement de
paramètre.

N.B. Deux tels chemins peuvent être représentés pour des physiciens par deux
points décrivant la même trajectoire mais avec des lois de mouvements différentes.

2.4. Si on fait un changement de paramètre vérifiant $t'(u) < 0$
$t(a_1) = b$, $t(b_1) = a$, on dit qu'on change l'orientation de $\gamma$ . Si $\gamma_1$
désigne le nouveau chemin on a alors $\int_{\gamma_1} \omega = - \int_\gamma \omega$ .

\uline{Définition 2.5.}

On appelle chemin différentiable par morceaux une application
continue $\gamma : [a,b] \to \mathbb{R}^2$ telle qu'il existe une subdivision de l'intervalle
$[a,b]$ $a = t_0 < t_1 < t_2 \dots < t_n = b$ de façon que les restrictions $\gamma$ aux
intervalles $[t_i, t_{i+1}]$ , $(i = 0,1,\dots,n-1)$, soient continûment différentiables
(i.e. vérifient 2.1).
On posera $\int_\gamma \omega = \sum_{i=0}^{n-1} \int_{\gamma_i} \omega$, et on vérifiera sans peine que $\int_\gamma \omega$ ne dépend
pas de la subdivision choisie.

3. \uline{Primitives d'une forme différentiable}

\uline{Lemme 3.1.} Soit $U$ un ouvert connexe du plan. Quels que soient les points
$a \in U$ et $b \in U$, il existe un chemin $\gamma$ différentiable par morceaux, contenu
dans $U$ ayant a pour origine et $b$ pour extrémité.

\margin{165}
